\documentclass[aspectratio=169]{beamer}
\usetheme{Madrid}
\usecolortheme{seahorse}
\usepackage[utf8]{inputenc}
\usepackage[T1]{fontenc}
\usepackage[polish]{babel}
\usepackage{booktabs}
\usepackage{graphicx}
\usepackage{amsmath}
\usepackage{array}
\setbeamertemplate{navigation symbols}{}
\setbeamertemplate{caption}[numbered]
\setbeamertemplate{footline}{\leavevmode\hbox{\begin{beamercolorbox}[wd=.86\paperwidth,ht=2.3ex,dp=1ex,leftskip=1em]{author in head/foot}Szkoła Główna Handlowa | Projekt zaliczeniowy ACT | czerwiec 2026\end{beamercolorbox}\begin{beamercolorbox}[wd=.14\paperwidth,ht=2.3ex,dp=1ex,center]{date in head/foot}\insertframenumber/\inserttotalframenumber\end{beamercolorbox}}}
\newcommand{\slidefig}[1]{\includegraphics[width=.88\textwidth,height=.62\textheight,keepaspectratio]{#1}}
\newcommand{\widefig}[1]{\includegraphics[width=.94\textwidth,height=.66\textheight,keepaspectratio]{#1}}
\newcommand{\slidetable}[1]{\begingroup\scriptsize\setlength{\tabcolsep}{2.4pt}\renewcommand{\arraystretch}{0.98}\input{#1}\endgroup}
\graphicspath{{../Wykresy/}}
\title{Analiza czasu trwania terapii i leczenia chorób serca}
\author{Michał Zwierzewicz \and Jakub Mikołajczak \and Patryk Piasecki}
\date{czerwiec 2026}
\begin{document}
\begin{frame}\titlepage\end{frame}
\begin{frame}{Agenda}
\begin{enumerate}
\item Cel badania i opis danych
\item Zmienna czasu, zdarzenie i cenzurowanie
\item Analiza nieparametryczna
\item Modele parametryczne AFT
\item Rozszerzenie bayesowskie
\item Model Coxa i diagnostyka
\item Porównanie modeli i wnioski
\end{enumerate}
\end{frame}
\begin{frame}{Cel badania}
\begin{itemize}
\item Zbiór zawiera 13 zmiennych źródłowych i 299 obserwacji pacjentów z niewydolnością serca.
\item Modelowana jest zmienna czasu \texttt{time}, a status zdarzenia określa \texttt{DEATH\_EVENT}.
\item \texttt{DEATH\_EVENT=1} oznacza zgon, a \texttt{DEATH\_EVENT=0} brak zgonu do końca obserwacji.
\item Brak zgonu do końca obserwacji traktowany jest jako cenzurowanie prawostronne.
\item Celem projektu jest porównanie nieparametrycznych, parametrycznych i semiparametrycznych metod analizy czasu trwania oraz ocena ich założeń.
\end{itemize}
\end{frame}
\begin{frame}{Zbiór danych}
\begin{columns}[T,totalwidth=\textwidth]
\begin{column}{.48\textwidth}
\begin{itemize}
\item Obserwacje: 299
\item Zgony: 96
\item Cenzurowanie prawostronne: 203
\item Zmienna czasu: \texttt{time}
\item Zmienna zdarzenia: \texttt{DEATH\_EVENT}
\end{itemize}
\end{column}
\begin{column}{.50\textwidth}
\begin{itemize}
\item Zmienne kliniczne: wiek, płeć, anemia, cukrzyca, nadciśnienie, palenie.
\item Zmienne laboratoryjne: frakcja wyrzutowa, kreatynina, sód, płytki krwi, CPK.
\item Transformacja CPK: \texttt{log\_cpk = log(CPK+1)}.
\end{itemize}
\end{column}
\end{columns}
\end{frame}
\begin{frame}{Zmienne definicyjne i binarne}
\begin{tabular}{p{.30\textwidth}p{.60\textwidth}}
\toprule
Zmienna & Znaczenie \\
\midrule
\texttt{DEATH\_EVENT} & 1: zgon; 0: obserwacja cenzurowana prawostronnie \\
\texttt{anaemia} & rozpoznanie anemii \\
\texttt{diabetes} & rozpoznanie cukrzycy \\
\texttt{high\_blood\_pressure} & nadciśnienie tętnicze \\
\texttt{sex} & płeć pacjenta; 1: mężczyzna, 0: kobieta \\
\texttt{smoking} & palenie tytoniu \\
\bottomrule
\end{tabular}
\end{frame}
\begin{frame}{Zmienne ilościowe i transformacje}
\begin{tabular}{p{.30\textwidth}p{.60\textwidth}}
\toprule
Zmienna & Znaczenie \\
\midrule
\texttt{time} & czas obserwacji w dniach \\
\texttt{age} & wiek pacjenta \\
\texttt{ejection\_fraction} & frakcja wyrzutowa serca \\
\texttt{serum\_creatinine} & poziom kreatyniny w surowicy \\
\texttt{serum\_sodium} & poziom sodu w surowicy \\
\texttt{platelets} & liczba płytek krwi \\
\texttt{creatinine\_phosphokinase} & poziom CPK; w modelach użyto \texttt{log\_cpk} \\
\bottomrule
\end{tabular}
\end{frame}
\begin{frame}{Statystyki opisowe zmiennych ilościowych}\slidetable{../Tables/01_statystyki_opisowe.tex}\end{frame}
\begin{frame}{Zmiennie binarne i zbinaryzowane}\slidetable{../Tables/01b_zmienne_binarne.tex}\end{frame}
\begin{frame}{Struktura zdarzeń i cenzurowania}
\begin{columns}[T,totalwidth=\textwidth]
\begin{column}{.42\textwidth}\slidetable{../Tables/02_struktura_cenzurowania.tex}\end{column}
\begin{column}{.56\textwidth}\centering\slidefig{02_struktura_cenzurowania.pdf}\end{column}
\end{columns}
{\scriptsize W zbiorze odnotowano 96 zgonów oraz 203 obserwacji cenzurowanych prawostronnie. Udział cenzurowania wpływa na precyzję estymacji funkcji przeżycia, szczególnie w końcowym zakresie czasu obserwacji.}
\end{frame}
\begin{frame}{Analiza nieparametryczna}
Analiza nieparametryczna stanowi punkt odniesienia dla dalszego modelowania. Pozwala ocenić empiryczny przebieg funkcji przeżycia, skumulowanego hazardu oraz różnice między wybranymi podgrupami bez narzucania konkretnego rozkładu czasu przeżycia.
\end{frame}
\begin{frame}{Estymator Kaplana-Meiera}
\[
\widehat{S}(t)=\prod_{t_i \le t}\left(1-\frac{d_i}{n_i}\right)
\]
\begin{itemize}
\item $d_i$ oznacza liczbę zdarzeń w czasie $t_i$.
\item $n_i$ oznacza liczbę osób w zbiorze ryzyka tuż przed $t_i$.
\item Obserwacje cenzurowane zmniejszają zbiór ryzyka, ale nie powodują skoku funkcji przeżycia.
\end{itemize}
\end{frame}
\begin{frame}{Kaplan-Meier dla całej próby}\centering\widefig{03_kaplan_meier.pdf}\\[-.3em]{\scriptsize Zacieniowany obszar oznacza przedział ufności dla estymowanej funkcji przeżycia.}\end{frame}
\begin{frame}{Tablica aktuarialna: wyniki}\slidetable{../Tables/04_tablica_aktuarialna.tex}\end{frame}
\begin{frame}{Tablica aktuarialna: wykres}\centering\widefig{05_tablica_aktuarialna.pdf}\end{frame}
\begin{frame}{Skumulowany hazard Nelsona-Aalena}
\[
\widehat{H}(t)=\sum_{t_i \le t}\frac{d_i}{n_i}
\]
Skumulowany hazard rośnie wraz z narastaniem ryzyka zdarzenia wśród pacjentów pozostających w obserwacji.
\end{frame}
\begin{frame}{Nelson-Aalen: wynik}\centering\widefig{04_nelson_aalen.pdf}\end{frame}
\begin{frame}{Nieparametryczna funkcja hazardu}
\centering
\widefig{04b_hazard_smoothed_bw50.pdf}\\[-.3em]
{\scriptsize Wykres przedstawia wygładzoną funkcję hazardu Nelsona-Aalena dla pasma $bandwidth=50$. Krzywa pokazuje bieżącą intensywność zgonów w czasie, a zacieniowany obszar oznacza 95\% przedział ufności.}
\end{frame}
\begin{frame}{Kaplan-Meier dla podgrup I}\centering\widefig{06_km_podgrupy_1.pdf}\end{frame}
\begin{frame}{Kaplan-Meier dla podgrup II}\centering\widefig{06_km_podgrupy_2.pdf}\end{frame}
\begin{frame}{Kaplan-Meier dla podgrup III}\centering\widefig{06_km_podgrupy_3.pdf}\end{frame}
\begin{frame}{Kaplan-Meier dla podgrup IV}\centering\widefig{06_km_podgrupy_4.pdf}\end{frame}
\begin{frame}{Testy log-rank i Wilcoxona}
\[
H_0:S_0(t)=S_1(t)\ \text{dla wszystkich}\ t,\qquad
H_1:S_0(t)\neq S_1(t)\ \text{dla co najmniej jednego}\ t.
\]
Test log-rank nadaje podobną wagę zdarzeniom w czasie, a test Wilcoxona silniej waży wcześniejsze zdarzenia.
\end{frame}
\begin{frame}{Wyniki testów nieparametrycznych}\slidetable{../Tables/05_testy_nieparametryczne.tex}\end{frame}
\begin{frame}{Wnioski z części nieparametrycznej}
Estymator Kaplana-Meiera pokazuje empiryczny spadek prawdopodobieństwa przeżycia bez narzucania rozkładu czasu trwania, natomiast estymator Nelsona-Aalena uzupełnia analizę o obraz narastania skumulowanego hazardu. Obie metody stanowią punkt odniesienia dla modeli parametrycznych i modelu Coxa.
\end{frame}
\begin{frame}{Modele parametryczne AFT}
Modele parametryczne opisują rozkład czasu przeżycia przez określoną rodzinę rozkładów. Ich zaletą jest możliwość estymacji funkcji przeżycia, hazardu i scenariuszy predykcyjnych dla profili pacjentów. Ich ograniczeniem jest wrażliwość na błędną specyfikację rozkładu czasu przeżycia.
\end{frame}
\begin{frame}{Model Weibull AFT}
\[
\log(T)=\beta_0+\beta^TX+\sigma\varepsilon,\qquad
\lambda(X)=\exp(\beta_0+\beta^TX)
\]
\[
\rho>1:\ \text{hazard rosnący},\quad \rho=1:\ \text{hazard stały},\quad \rho<1:\ \text{hazard malejący}
\]
W modelu AFT $\exp(\beta)$ jest współczynnikiem czasu, a nie hazard ratio.
\end{frame}
\begin{frame}{Porównanie modeli parametrycznych}\slidetable{../Tables/07_porownanie_modeli_parametrycznych.tex}\end{frame}
\begin{frame}{Interpretacja porównania parametrycznego}
Najniższe wartości kryteriów informacyjnych uzyskał model Weibull AFT. Model wykładniczy bez kowariant estymuje jeden parametr intensywności zdarzenia, ponieważ nie zawiera zmiennych objaśniających ani dodatkowego parametru kształtu.
\end{frame}
\begin{frame}{Ocena dopasowania rozkładu -- probability plots}
\centering
\widefig{08f_probability_plots.pdf}\\[-.3em]
{\scriptsize Probability plots służą do wizualnej oceny zgodności empirycznego rozkładu czasu przeżycia z rozkładami parametrycznymi. Im bliżej punktów do linii referencyjnej, tym lepsze dopasowanie rozkładu w sensie marginalnym.}
\end{frame}
\begin{frame}{Wybór rozkładu: Weibull AFT}
\begin{itemize}
\item Probability plots wskazują dobre dopasowanie graficzne trzech rozkładów parametrycznych; wartości $R^2$ przekraczają 0{,}93.
\item Właściwym kryterium wyboru modelu regresyjnego jest AIC dla pełnego modelu z kowariantami.
\item Pełny model Weibull AFT ma najniższe AIC w tej grupie; przewaga względem modelu log-normalnego wynosi około $\Delta AIC=5{,}7$.
\item Log-logistic AFT pozostaje modelem konkurencyjnym, natomiast log-normal AFT ma dobre dopasowanie graficzne, ale słabsze AIC.
\item Oszacowany parametr kształtu $\hat{\rho}\approx0{,}96$ jest bliski 1, co wskazuje na hazard zbliżony do stałego i klinicznie interpretowalną strukturę modelu.
\end{itemize}
\end{frame}
\begin{frame}{Parametry pełnego modelu Weibull AFT}\slidetable{../Tables/08_parametry_najlepszego_aft.tex}\end{frame}
\begin{frame}{Model Weibull AFT zredukowany}\slidetable{../Tables/08b_parametry_aft_zredukowany.tex}\end{frame}
\begin{frame}{Model pełny i zredukowany}\slidetable{../Tables/08c_aft_pelny_vs_zredukowany.tex}\end{frame}
\begin{frame}{Profile przeżycia w modelu AFT}\centering\widefig{08c_weibull_profile_survival.pdf}\\[-.3em]{\scriptsize Profil wysokiego ryzyka charakteryzuje się szybszym spadkiem funkcji przeżycia w modelu Weibull AFT zredukowanym.}\end{frame}
\begin{frame}{Hazard w modelu AFT}\centering\widefig{08e_weibull_profile_hazard.pdf}\\[-.3em]{\scriptsize Profile ryzyka różnią się poziomem hazardu, przy podobnym kształcie przebiegu wynikającym ze wspólnego parametru rozkładu Weibulla.}\end{frame}
\begin{frame}{Diagnostyka parametryczna}\centering\widefig{08b_probability_diagnostics.pdf}\end{frame}
\begin{frame}{Ograniczenie liczby parametrów}
Liczba zdarzeń ogranicza liczbę parametrów, które można stabilnie interpretować. Zbyt rozbudowany model może prowadzić do niestabilnych oszacowań, szerokich przedziałów ufności i nadmiernego dopasowania do próby.
\end{frame}
\begin{frame}{Rozszerzenie bayesowskie}
Bayesowski model Weibulla został oszacowany w uproszczonej wersji dla trzech kluczowych zmiennych: wieku, frakcji wyrzutowej i kreatyniny. Model ma charakter demonstracyjny i służy ocenie kierunków efektów oraz niepewności posteriori.
\end{frame}
\begin{frame}{Bayesowski model Weibulla}
\[
T_i\sim Weibull(\rho,\lambda_i),\qquad \log\lambda_i=\beta_0+\beta^TX_i
\]
\[
\beta_j\sim N(0,2.5^2),\qquad \log\rho\sim N(0,1)
\]
Cenzurowanie uwzględniono przez wkład $S(t_i|X_i)$ dla obserwacji bez zdarzenia.
\end{frame}
\begin{frame}{Posterior modelu bayesowskiego}\slidetable{../Tables/08d_bayes_posterior.tex}\end{frame}
\begin{frame}{Traceploty MCMC}\centering\widefig{08d_bayes_traceplots_refreshed.pdf}\\[-.3em]{\scriptsize Średnia akceptacja MCMC wyniosła 0{,}5339, czyli przyjęto około 53{,}4\% proponowanych przejść algorytmu Metropolisa-Hastingsa. Taki poziom oznacza, że łańcuch eksploruje posterior bez nadmiernego odrzucania propozycji.}\end{frame}
\begin{frame}{Model Coxa}
Model Coxa umożliwia wielowymiarową analizę wpływu zmiennych klinicznych na hazard zgonu bez konieczności parametryzacji hazardu bazowego. Ceną tej elastyczności jest konieczność weryfikacji założenia proporcjonalności hazardów.
\end{frame}
\begin{frame}{Model Coxa: idea}
\[
h(t|X)=h_0(t)\exp(\beta^TX)
\]
$\exp(\beta_j)$ to hazard ratio dla jednostkowej zmiany zmiennej $X_j$, przy pozostałych zmiennych stałych.
\end{frame}
\begin{frame}{Wyniki modelu Coxa}\slidetable{../Tables/09_model_coxa.tex}\end{frame}
\begin{frame}{Współczynniki modelu Coxa}\centering\widefig{09_cox_wspolczynniki.pdf}\\[-.3em]{\scriptsize Wartości dodatnie oznaczają wzrost hazardu, a ujemne jego spadek; poziome odcinki oznaczają przedziały ufności.}\end{frame}
\begin{frame}{Założenie proporcjonalnych hazardów}
Poprawna kolejność analizy obejmuje estymację modelu Coxa, diagnostykę założenia PH, ewentualną modyfikację modelu i ponowną estymację. Naruszenie PH oznacza, że stałe hazard ratio może być interpretowane jedynie jako efekt uśredniony.
\end{frame}
\begin{frame}{Test proporcjonalności hazardów}\slidetable{../Tables/10_test_ph_coxa.tex}\vspace{-0.4em}{\scriptsize Test wskazuje na naruszenie założenia PH dla frakcji wyrzutowej na poziomie istotności 5\%.}\end{frame}
\begin{frame}{Reszty Schoenfelda I}\centering\widefig{10_reszty_schoenfelda_1.pdf}\end{frame}
\begin{frame}{Reszty Schoenfelda II}\centering\widefig{10_reszty_schoenfelda_2.pdf}\end{frame}
\begin{frame}{Reszty Schoenfelda III}\centering\widefig{10_reszty_schoenfelda_3.pdf}\end{frame}
\begin{frame}{Reszty Schoenfelda IV}\centering\widefig{10_reszty_schoenfelda_4.pdf}\end{frame}
\begin{frame}{Reszty Schoenfelda V}\centering\widefig{10_reszty_schoenfelda_5.pdf}\end{frame}
\begin{frame}{Model Coxa z interakcją czasową}\slidetable{../Tables/10b_cox_interakcja_czasowa.tex}\end{frame}
\begin{frame}{Wyniki modelu Coxa z interakcją czasową}\slidetable{../Tables/10c_cox_interakcja_wyniki.tex}\end{frame}
\begin{frame}{Test PH po dodaniu interakcji czasowej}\slidetable{../Tables/10d_test_ph_cox_interakcja.tex}\vspace{-0.4em}{\scriptsize Ponowny test służy ocenie, czy wprowadzenie efektu zależnego od czasu osłabiło sygnał naruszenia założenia proporcjonalności hazardów.}\end{frame}
\begin{frame}{Obserwacje wpływowe}\slidetable{../Tables/11_obserwacje_wplywowe.tex}\vspace{-0.4em}{\scriptsize Miary delta-beta pokazują, jak zmieniłyby się współczynniki po usunięciu danej obserwacji.}\end{frame}
\begin{frame}{Predykcja funkcji przeżycia z modelu Coxa}\centering\widefig{11_cox_profile_survival_three_risk.pdf}\\[-.3em]{\scriptsize Trzy klasy ryzyka pokazują praktyczną interpretację modelu Coxa: wraz z pogarszaniem profilu klinicznego przewidywana funkcja przeżycia opada coraz szybciej.}\end{frame}
\begin{frame}{Predykcja funkcji hazardu z modelu Coxa}\centering\widefig{11b_cox_profile_hazard_three_risk.pdf}\\[-.3em]{\scriptsize Funkcja hazardu pokazuje bieżącą intensywność zdarzenia dla profili ryzyka. W modelu Coxa profile różnią się poziomem hazardu przez czynnik $\exp(\beta^T X)$, przy wspólnym hazardzie bazowym.}\end{frame}
\begin{frame}{Porównanie klas modeli}\slidetable{../Tables/12_porownanie_syntetyczne_modeli.tex}\end{frame}
\begin{frame}{Ograniczenia analizy}
\begin{itemize}
\item 299 obserwacji i 96 zdarzeń ogranicza liczbę stabilnie interpretowanych parametrów.
\item Duży udział cenzurowania obniża precyzję estymacji w końcowym zakresie czasu.
\item Założenie niezależnego cenzurowania nie jest w pełni weryfikowalne.
\item Modele parametryczne zależą od wyboru rozkładu czasu przeżycia.
\item Dla frakcji wyrzutowej występuje sygnał naruszenia założenia PH.
\item Model bayesowski ma charakter demonstracyjny i obejmuje ograniczony zestaw zmiennych.
\end{itemize}
\end{frame}
\begin{frame}{Elementy kontroli jakości analizy}
\begin{itemize}
\item Spójnie zdefiniowano czas, zdarzenie i cenzurowanie.
\item Tabele i wykresy prezentacji są generowane w oficjalnych notebookach.
\item Model AFT oceniono w wersji pełnej i zredukowanej.
\item Model Coxa uzupełniono diagnostyką PH, resztami Schoenfelda i wariantem z efektem czasowym.
\item Rozszerzenie bayesowskie obejmuje estymację MCMC i diagnostykę traceplotów.
\end{itemize}
\end{frame}
\begin{frame}{Podsumowanie}
Analiza przeżycia została przeprowadzona w trzech uzupełniających się podejściach. Metody nieparametryczne opisały empiryczny przebieg przeżycia, model Weibull AFT pozwolił na parametryczną predykcję czasu przeżycia, a model Coxa umożliwił interpretację wpływu zmiennych przez hazard ratio. Wyniki wskazują na szczególne znaczenie wieku, frakcji wyrzutowej i kreatyniny, przy konieczności ostrożnej interpretacji założenia proporcjonalności hazardów dla frakcji wyrzutowej.
\end{frame}
\end{document}
